\chapter{Análise do Sistema}
\label{cap:analise}

\orientacao{Este capítulo transforma a compreensão do contexto e das necessidades em uma visão estruturada do sistema a ser desenvolvido. Sua função é definir o que o sistema deve fazer, quem interage com ele, quais conceitos compõem seu domínio e quais requisitos e restrições devem ser considerados, sem ainda detalhar sua implementação.}

\section{Visão Geral do Sistema}
\orientacao{Apresentar uma visão macroscópica da solução proposta, delimitando o sistema e sua finalidade antes do detalhamento de suas funcionalidades e estruturas.}

\section{Atores Envolvidos}
\orientacao{Identificar os atores que interagem diretamente com o sistema e caracterizar seus respectivos papéis.}

\section{Casos de Uso}
\orientacao{Identificar e descrever as funcionalidades do sistema sob a perspectiva dos atores, incluindo os comportamentos esperados e os fluxos relevantes.}

\diagrama{Diagrama de casos de uso}{casos-de-uso}

\section{Modelo Conceitual}
\orientacao{Representar os principais conceitos do domínio, suas características e relacionamentos, sem incorporar decisões específicas de implementação.}

\diagrama{Diagrama de classes conceitual}{classes-conceitual}

\section{Requisitos Funcionais}
\orientacao{Especificar de forma objetiva os serviços e comportamentos que o sistema deverá oferecer para atender às necessidades identificadas.}

\section{Requisitos Não Funcionais}
\orientacao{Especificar as características, qualidades e restrições que o sistema deverá atender, além das funcionalidades propriamente ditas.}

\section{Restrições do Sistema}
\orientacao{Registrar limitações impostas ao desenvolvimento ou à operação do sistema, sejam elas técnicas, institucionais, legais, orçamentárias ou decorrentes do contexto do projeto.}
